
\section*{第二章\ 新增頁碼與編輯不同類別頁碼}
\addcontentsline{toc}{section}{第二章\ 新增頁碼與編輯不同類別頁碼}

\subsection*{2.1\ 分隔不同頁面以使用不同種類頁碼}
\addcontentsline{toc}{subsection}{2.1\ 分隔不同頁面以使用不同種類頁碼}

\subsubsection*{2.1.1\ 分隔頁面}
\addcontentsline{toc}{subsection}{2.1.1\ 分隔頁面}
為了新增不同種類的頁碼，我們需要先分隔頁面，首先先點擊摘要頁的標題開頭，然後點擊「版面配置」中的「分隔符號」，選擇「下一頁」，如此一來封面和摘要頁的頁碼格式就可以分開來了，同樣的事情在內文的第一頁也須做一次，才能在摘要和圖表目錄間的頁面使用羅馬數字格式的頁碼。
 
\subsubsection*{2.1.2\ 新增頁碼}
\addcontentsline{toc}{subsection}{2.1.2\ 新增頁碼}
接著就是新增頁碼，點選「插入」-「頁碼」，選擇頁碼格式，依據所在頁面選擇符合格式的頁碼種類後，點選「頁面底端」的「純數字2」，就可以新增適合的頁碼到頁面底部。
\clearpage